% Preámbulo del informe de situación de f1-data-platform (compilar con XeLaTeX o latexmk -xelatex)
\usepackage[a4paper,margin=2.4cm,headheight=14pt]{geometry}
\usepackage{fontspec}
\usepackage[spanish,es-noshorthands,es-tabla]{babel}
\setmainfont{Libertinus Serif}
\setsansfont{Fira Sans}[Scale=MatchLowercase]
\setmonofont{Fira Mono}[Scale=0.86]
\usepackage{microtype}
\usepackage[dvipsnames,table]{xcolor}
\usepackage{graphicx}
\usepackage{booktabs,tabularx,longtable,array,multirow}
\usepackage{enumitem}
\setlist{itemsep=2pt,topsep=4pt}
\usepackage{csquotes}
\usepackage{fancyhdr}
\usepackage{titlesec}
\usepackage{caption}
\captionsetup{font=small,labelfont={bf,sf}}
\usepackage{listings}
\usepackage{tcolorbox}
\tcbuselibrary{skins,breakable}
\usepackage{tikz}
\usetikzlibrary{positioning,arrows.meta,shapes.geometric,shapes.multipart,fit,backgrounds,calc,matrix,chains}
\usepackage{pgfplots}
\pgfplotsset{compat=1.18}
\usepackage{amssymb,pifont}
\usepackage{fvextra}
\usepackage[colorlinks,urlcolor=oiblue!75!black,linkcolor=ink,citecolor=ink]{hyperref}
\usepackage{bookmark}

% Colores del proyecto (paleta Okabe-Ito + rojo F1)
\definecolor{f1red}{HTML}{C8102E}
\definecolor{ink}{HTML}{1F2330}
\definecolor{muted}{HTML}{5B6170}
\definecolor{bronze}{HTML}{A0652A}
\definecolor{silver}{HTML}{7C8691}
\definecolor{gold}{HTML}{B8901A}
\definecolor{oiblue}{HTML}{0072B2}
\definecolor{oiorange}{HTML}{E69F00}
\definecolor{oigreen}{HTML}{009E73}
\definecolor{oisky}{HTML}{56B4E9}
\definecolor{oiverm}{HTML}{D55E00}
\definecolor{oipurple}{HTML}{CC79A7}
\definecolor{codebg}{HTML}{F5F6F8}

\titleformat{\chapter}[display]{\sffamily\bfseries\huge\color{ink}}{\color{f1red}\large\MakeUppercase{\chaptertitlename}\ \thechapter}{4pt}{}
\titleformat{\section}{\sffamily\bfseries\Large\color{ink}}{\color{f1red}\thesection}{0.8em}{}
\titleformat{\subsection}{\sffamily\bfseries\large\color{ink}}{\color{muted}\thesubsection}{0.8em}{}
\titleformat{\subsubsection}{\sffamily\bfseries\normalsize\color{ink}}{}{0em}{}
\setcounter{secnumdepth}{2}
\setcounter{tocdepth}{2}

\pagestyle{fancy}
\fancyhf{}
\fancyhead[L]{\sffamily\small\color{muted}\leftmark}
\fancyhead[R]{\sffamily\small\color{muted}f1-data-platform · Informe de situación}
\fancyfoot[C]{\sffamily\small\thepage}
\renewcommand{\headrulewidth}{0.3pt}

\lstset{
  basicstyle=\ttfamily\small,
  backgroundcolor=\color{codebg},
  frame=none,
  breaklines=true,
  columns=fullflexible,
  keepspaces=true,
  showstringspaces=false,
  xleftmargin=6pt,
  framexleftmargin=6pt,
  aboveskip=8pt,
  belowskip=8pt,
  keywordstyle=\color{oiblue}\bfseries,
  commentstyle=\color{muted}\itshape,
  stringstyle=\color{oiverm},
  literate={á}{{\'a}}1 {é}{{\'e}}1 {í}{{\'i}}1 {ó}{{\'o}}1 {ú}{{\'u}}1 {ñ}{{\~n}}1 {Á}{{\'A}}1 {É}{{\'E}}1 {Í}{{\'I}}1 {Ó}{{\'O}}1 {Ú}{{\'U}}1 {Ñ}{{\~N}}1 {ü}{{\"u}}1 {«}{{\guillemotleft}}1 {»}{{\guillemotright}}1 {→}{{$\to$}}1 {×}{{$\times$}}1
}

% Cajas de resumen, recomendación y aviso
\newtcolorbox{resumen}[1][Resumen]{enhanced,breakable,colback=codebg,colframe=ink,boxrule=0.4pt,arc=2pt,
  fonttitle=\sffamily\bfseries,title={#1},left=8pt,right=8pt}
\newtcolorbox{recomendacion}[1][Recomendación]{enhanced,breakable,colback=oigreen!6,colframe=oigreen!70!black,boxrule=0.4pt,arc=2pt,
  fonttitle=\sffamily\bfseries,title={#1},left=8pt,right=8pt}
\newtcolorbox{aviso}[1][Atención]{enhanced,breakable,colback=oiorange!8,colframe=oiorange!80!black,boxrule=0.4pt,arc=2pt,
  fonttitle=\sffamily\bfseries,title={#1},left=8pt,right=8pt}

% Estilos TikZ comunes
\tikzset{
  caja/.style={draw=ink,rounded corners=2pt,align=center,font=\sffamily\small,inner sep=5pt,fill=white},
  capa/.style={draw=none,rounded corners=4pt,inner sep=8pt},
  flecha/.style={-{Stealth[length=2.4mm]},thick,color=ink!80},
  flechad/.style={-{Stealth[length=2.4mm]},thick,dashed,color=muted},
  etiqueta/.style={font=\sffamily\scriptsize,color=muted},
  tabla/.style={rectangle split,rectangle split parts=2,draw=ink,rounded corners=1pt,align=left,font=\sffamily\scriptsize,
    rectangle split part fill={#1,white},inner sep=3pt,text width=3.3cm},
}

\newcommand{\code}[1]{\texttt{\small\color{ink!85!oiblue}#1}}
\newcommand{\urltexto}[1]{\texttt{\footnotesize #1}}
\newcommand{\si}{\textcolor{oigreen!80!black}{\ding{51}}}
\newcommand{\no}{\textcolor{oiverm}{\ding{55}}}
\newcommand{\parcial}{\textcolor{oiorange}{\ding{108}}}
\newcommand{\nada}{\textcolor{muted}{\ding{109}}}

% Bloques de código de las notas (fvextra corta las líneas largas)
\DefineVerbatimEnvironment{codigo}{Verbatim}{breaklines,breakanywhere,fontsize=\footnotesize,
  frame=leftline,framerule=1.2pt,rulecolor=\color{oiblue!50},xleftmargin=8pt,framesep=6pt}
\DefineVerbatimEnvironment{diagramatexto}{Verbatim}{breaklines,breakanywhere,fontsize=\scriptsize,
  frame=single,rulecolor=\color{muted!50},framesep=5pt,label={\sffamily\scriptsize descripción del diagrama}}
\newenvironment{nota}{\begin{tcolorbox}[enhanced,breakable,colback=codebg,colframe=codebg,
  borderline west={2pt}{0pt}{muted},arc=0pt,left=8pt,right=6pt,top=3pt,bottom=3pt,fontupper=\small]}{\end{tcolorbox}}
\newcommand{\fich}[1]{\texttt{\small\color{oiblue!70!black}#1}}
\newcommand{\prio}[1]{\textsf{\small\bfseries #1}}
\setlength{\emergencystretch}{3em}
\counterwithout{figure}{chapter}
